\documentclass[graybox,envcountchap,sectrefs]{svmono}

\usepackage[spanish]{babel}
\usepackage[utf8]{inputenc}
\usepackage{amsmath, amsfonts, amssymb}
\usepackage{minted}
\usepackage{xcolor}
\usepackage{tcolorbox}
\tcbuselibrary{minted}
\definecolor{latteBase}{HTML}{EFF1F5}
\newtcblisting{code}{
    listing engine=minted,
    minted language=latex,
   % minted options={fontsize=\small},
    colback=latteBase,
    %colback=black!20,
    colframe=black!5,
    boxrule=0.5pt,
    left=5mm,
    listing only
}

\begin{document}

\author{José Alejandro Aburto}
\title{Cursillo \LaTeX}
\subtitle{Una introducción enfoncada en docencia en matemáticas}
\maketitle

\chapter*{Introducción}

\chapter{Estructura de un documento \LaTeX}
Un documento \LaTeX tiene la extensión \textbf{.tex}
y se compone de dos partes:
    \begin{enumerate}
        \item El preámbulo.
        \item El cuerpo.
    \end{enumerate}
    En el preámbulo indicamos el tipo de documento que estamos realizando. Por ejemplo, en las plantillas
    que utilizaremos usamos la clase \emph{amsart} que el formato para artículo de la AMS (American Mathematical Society).
    El código que indica esto es el siguiente:
        \begin{code}
            \documentclass{amsart}
        \end{code}
        También podemos agregar información como

\chapter{Ecuaciones}
\section{Ecuaciones más usadas}
\section{Matrices}
\section{Tablas}
\section{Arreglos y casos}

\chapter{Itemizaciones, Enumeraciones y Gráficos}
\section{Itemize}
\section{Enumerate}
\section{includegraphics}

\chapter{TikZ*}
\section{Gráfico de Funciones}
\section{tkz-euclide}

%agrega de bibliografia la documentacion de:
% tikz, tkz-euclide, latex, y de la libreria de tikz con la que se hacen los graficos
% % ojo que puede que necesite de --shell-scape para este ultimo capitulo. 
% Probablemente sea mejor dejarlo como opcional o removerlo
\end{document}

